\documentclass[12pt,a4paper,oneside]{article}
\usepackage[brazil]{babel}
\usepackage[utf8]{inputenc}
\usepackage[dvipsnames]{xcolor}
\usepackage{listings}
\usepackage{wrapfig}
\usepackage{olymp}

\setcounter{problem}{1}

\begin{document}

\begin{problem}{l33t}{l33t}{2}

\begingroup
\setlength{\intextsep}{0pt}%
%\setlength{\columnsep}{0pt}

O \textit{Leet} (também chamado de \textit{l33t} ou $1337$) é uma forma de escrita criada na internet em que letras são substituídas por números e símbolos visualmente semelhantes.
Por exemplo, o número $4$ pode representar a letra `\texttt{A}', $3$ o `\texttt{E}', $1$ o `\texttt{I}' ou `\texttt{L}', e $0$ o `\texttt{O}'.

Essa técnica surgiu nas primeiras comunidades online, especialmente entre hackers e entusiastas de tecnologia,
para evitar filtros automáticos, criar uma linguagem interna de grupo e demonstrar conhecimento técnico (já que ``elite'' virou ``leet'' e depois ``1337'' (similar a \texttt{LEET} de cabeça para baixo).
Com o tempo, passou a ser usada também em jogos online, apelidos e memes, muitas vezes de forma irônica ou estilizada.

Hoje, o \textit{leet} é mais um elemento cultural da internet do que uma necessidade prática, mas ainda aparece em nomes de usuário, desafios e exercícios de programação.

Nesta questão, considere somente as seguintes substituições, que são as mais comuns:

\begin{itemize}
\item `\texttt{E}' $\Rightarrow 3$
\item `\texttt{A}' $\Rightarrow 4$
\item `\texttt{I}' $\Rightarrow 1$
\item `\texttt{O}' $\Rightarrow 0$
\item `\texttt{S}' $\Rightarrow 5$
\item `\texttt{T}' $\Rightarrow 7$
\item `\texttt{L}' $\Rightarrow 1$
\item `\texttt{B}' $\Rightarrow 8$
\item `\texttt{G}' $\Rightarrow 6$
\item `\texttt{Z}' $\Rightarrow 2$
\end{itemize}

% \Notes

\InputFile

A entrada começa com um linha contendo um número inteiro $N$, seguida por $N$ palavras, uma em cada linha.

Restrições: $1 \leq N \leq 10$; as palavras contém apenas letras maiúsculas e nenhuma possui mais de $10$ caracteres.

\endgroup

\OutputFile

Para cada palavra da entrada, seu programa deve escrever a correspondente após as substituições.

% \Notes

\Example

\begin{example}
\exmp{
2
HACKER
BRASIL
}{
H4CK3R
8R4511
}%
\end{example}

\begin{example}
\exmp{
3
GOOGLE
PASSWORD
UFV
}{
600613
P455W0RD
UFV
}%
\end{example}

%Comentários sobre o exemplo, se necessário.

\end{problem}

\end{document}